\section{Fixed-geometry comparison}\label{sec:geometry}

The center magnetization isolates the microscopic cost of one spin
insertion.  We compare boxes by placing the same finite nest of annuli
around each corresponding source.  An attachment calculation separates
the ratio of a nest from that of its exterior; the nest ratios then
cancel, leaving exterior energies at fixed geometry.  The first comparison
is between square sizes, and the second concerns the free-box thermodynamic
correlation of two insertions.  All continuum limits used here keep the
holes and every separating region at fixed positive size.

\begin{proposition}[Comparison with center magnetization]
\label{geo:comparison}
For integer sequences \(n,n'\to\infty\) with \(n'/n\to d\in(0,\infty)\),
\begin{equation}\label{geo:center-ratio}
                  \frac{M_{n'}}{M_n}\longrightarrow d^{-1/8}.
\end{equation}
There is a constant \(c_{\rm shape}\in(0,\infty)\) such that
\begin{equation}\label{geo:bulk-ratio}
                  \frac{C_{b_c}(r)}{M_r^2}\longrightarrow c_{\rm shape}
                  \qquad(r\to\infty,\ r\in\Z).
\end{equation}
Here \(C_{b_c}\) is the thermodynamic limit in
\eqref{eq:critical-definition}; that limit precedes \(r\to\infty\).
In particular, \(M_n=n^{-1/8+o(1)}\).
\end{proposition}

\subsection{Separating the insertion from the exterior}

For a graph with the boundary and constant conventions of
Section~\ref{sec:heights}, write
\(\Phi_G^\tau=Z_G^\tau/Z_G^0\) for its bare sector ratio.
The fixed-domain sector input evaluates this ratio on a free exterior.
Let \(U\) be a bounded connected rectilinear domain with a simple outer
boundary and finitely many disjoint simple holes, all with dyadic vertices,
positive side lengths, and positive passages.  Give the height graph
free walls.  If \(\vartheta_0\) is a smooth closed one-form near
\(\overline U\) whose periods about the holes are prescribed multiples
of \(2\pi\), set
\[
 \mathcal M_U(\vartheta_0)
   =\inf_{g\in H^1(U)}\int_U|\vartheta_0+\nabla g|^2.
\]
Let \(\mathcal G_\nu(U)\) be its nearest-neighbor primary graph at mesh
\(\nu^{-1}\).  For integral lattice connections that are flat on the
retained graph away from the holes and have these periods,
the sector theorem of \cite[Theorem~3.1]{S} states that
\begin{equation}\label{geo:free-sector}
 \Phi_{\mathcal G_\nu(U)}^\tau
       \longrightarrow \exp\{-\mathcal M_U(\vartheta_0)/(2a)\}
       \qquad(\nu\to\infty),
\end{equation}
where \(a=8\pi\).  The limit holds along
every coherent bounded-step rounding of this fixed geometry.  The
energy depends only on the periods.  This statement concerns free
height walls; a pinned exterior will be treated below.

To connect \eqref{geo:free-sector} to a microscopic source, use the
annular construction of \cite[Sections~4 and~5]{S}.  Fix a small
\(\varepsilon>0\) and a rectilinear disk \(P\) with dyadic vertices
such that
\begin{equation}\label{geo:polygon}
 B(0,1-\varepsilon)\subset P\subset B(0,1+\varepsilon).
\end{equation}
Choose a sufficiently large power of four \(\lambda\), and an integer
\(\varrho_0\) aligned so that all relevant array divisions lie between
rows of dual vertices.  Put \(\varrho_i=\varrho_0\lambda^i\).
The \(i\)-th annulus has walls \(\varrho_{i-1}P\) and
\(\varrho_iP\).  Write \(Q_s=(-s,s)^2\) and
\(A_\lambda=\lambda P\setminus P^\circ\).  In the reference annulus
\(A_\lambda\), the guard is
\(\overline{Q_{2\sqrt\lambda}}\setminus Q_{\sqrt\lambda}\), and the
inner and outer end regions are \(A_\lambda\cap\overline{Q_4}\) and
\(A_\lambda\setminus Q_{\lambda/4}\), respectively.  Scale these
three regions by \(\varrho_{i-1}\), and call the guard \(\mathsf P_i\).
It is a full separating band, and each end contains every endpoint
in that annulus of a bond to be restored at its wall.  These prescribed
bands, with positive gaps, are the geometry used in the flat-pin
estimate below.

Choose one vertex in \(\mathsf P_i\).  A guard datum records every
height difference from that vertex.  In the annular quotient choose
its height to be zero, so the datum specifies every guard height.
The nearest-neighbor weight then factors between the two sides
conditionally on that datum.  Write \(\mu_i^\tau\) for the guard law
in the cut annulus, with its one common constant quotiented out.

A core is a connected finite block filling the first hole; it may
carry the source and has no guard data.
At an interface in this construction, let \(T^\tau(u)\) be the restored
partition sum divided by the cut partition sum at its incident guard
data \(u\),
and put \(d_T=T^0(0)\).  The restored sum includes the relative
integer translations of the joined components.  The same definition
applies to one joint interface joining several last annuli to a
connected exterior, with no direct bonds between distinct incident
annuli.  Its data are those in the incident annuli, while all exterior
heights are summed together inside that one kernel.
For a vertex region \(H\), let \([H]\) denote the event that all its
heights are equal, without prescribing their common value.  In a
centered cut annulus with guard \(\mathsf P_i\), put
\[
 \overline D_i(H,\mathsf P_i)
   =\log\frac{\Prob([H]\mid[\mathsf P_i])}{\Prob([H])}.
\]
The upper kernel bound of \cite[Lemma~5.1]{S}, its Equation~(8), holds
for all these interfaces, including a core interface.  For an
annulus--annulus or joint exterior interface, its Equation~(9) also
gives the centered average bound:
\begin{equation}\label{geo:kernel-bounds}
 0<\frac{T^\tau(u)}{d_T}\le1,
 \qquad
 \int\frac{T^0}{d_T}\prod_i\dd\mu_i^0
       \ge \exp\left\{-\sum_i\overline D_i(H_i,\mathsf P_i)\right\}.
\end{equation}
The sum and product use the incident annuli and the appropriate end
in each.  The first bound is pointwise in the data, and the second
is an average under independent centered cut annuli.  The centered
divisor \(d_T\) therefore works for the shifted kernel as well.
The flat-pin estimate of \cite[Proposition~4.1]{S}
makes each interaction on the right arbitrarily small in the mesh
limit by taking \(\lambda\) large at fixed \(P\).

Two consequences of \cite[Proposition~5.3 and Corollary~5.5]{S}
control how the averaged loss passes through a whole nest.
For any \(\xi\in(0,1/64)\), one can first take
\(\varepsilon\) small, then \(\lambda\) large, and finally the fixed
aligned threshold \(\varrho_0\) large, so that every interface between
successive annuli of these radii satisfies
\begin{equation}\label{geo:regular-loss}
 \int\left(1-\frac{T_i^\tau}{d_{T_i}}\right)
             \dd(\mu_i^\tau\otimes\mu_{i+1}^\tau)\le\xi.
\end{equation}
This holds in the centered sector and in either unit-period sector,
when restoration introduces no source.  It follows from the exact
sector identity (10) in that proposition and
\eqref{geo:free-sector}; it is the normalized estimate (11) there.
All late annuli obey the same bound because they refine the same
fixed geometry.  Fix one such core, containing the source in a unit
sector, and fix its connection once on the increasing family.  In a nest
of \(J\) annuli, the law of its last guard then has density at most \(2\)
with respect to
\(\mu_J^\tau\), for all sufficiently large \(J\).  This holds in
all three sectors and for identical lattice translates.  The only
lower assumption on the first core kernel is its strict positivity;
no uniform positive lower bound on that kernel is required.
This density bound compares the actual last-guard law, after all inner
blocks have been joined, with the cut law used in the averaged estimate.

It remains to attach the whole nests to their exterior.
Let \(\mathcal N_J\) be the bare sector ratio of one whole
free nest of depth \(J\), with its core and a unit source.  Integral
gauge changes and reversal of all heights show that the same
\(\mathcal N_J\) is obtained from any lattice translate and either
sign.  This is the common factor that will cancel in both comparisons.
Suppose \(k\) such nests, with disjoint closures, attach to a
connected exterior.  The graph of cut blocks, with an edge for each
interface, is a tree rooted at the exterior; there are no other block
adjacencies, in particular no direct bonds between distinct incident
last annuli.  Write \(\mu_{J,i}^\tau\) for the cut guard law in
the \(i\)-th incident last annulus, with its restricted connection.
In the following ratio, ``joined'' means the
exterior after joining only the \(k\) last annuli; the holes then have
the inner radii of those annuli.  Define
\begin{equation}\label{geo:joining-ratio}
 \mathfrak q^\tau_{\rm ext}
  =\frac{\Phi_{\rm joined}^\tau}
         {\Phi_{\rm ext}^\tau\prod_{i=1}^k\Phi_{{\rm last},i}^\tau}.
\end{equation}
Releasing those interfaces gives \(0<\mathfrak q^\tau_{\rm ext}\le1\).
Exact summation of the interface kernel, which is Equation~(10) of
\cite{S}, gives
\begin{equation}\label{geo:kernel-sector-identity}
 \int\frac{T^\tau_{\rm ext}}{d_{\rm ext}}
              \prod_i\dd\mu_{J,i}^\tau
 =\left(\int\frac{T^0_{\rm ext}}{d_{\rm ext}}
              \prod_i\dd\mu_{J,i}^0\right)\mathfrak q^\tau_{\rm ext}.
\end{equation}
Thus a sufficient attachment condition, along the lattice sequence
in question, is
\begin{equation}\label{geo:attachment-condition}
 \sum_i\overline D_i(H_i,\mathsf P_i)\le\delta+o(1),
 \qquad \log\mathfrak q^\tau_{\rm ext}\ge-\delta-o(1),
\end{equation}
with \(\delta>0\) arbitrarily small.  Indeed,
\eqref{geo:kernel-bounds} and \eqref{geo:kernel-sector-identity}
bound the shifted average loss by
\(1-e^{-2\delta-o(1)}\), and the centered one by
\(1-e^{-\delta-o(1)}\).  This estimate is averaged over the independent
cut-annulus laws, without any uniform lower bound at individual data.

The lower thermodynamic comparison below retains a hard pin in the
exterior.  We therefore check that the same reasoning applies when
that pin fixes the exterior constant and is kept in both the restored
and cut graphs.  The entire hard constraint remains in the exterior
partition sum and joint kernel; for the rim below, this includes every
equality \(h_x=0\) on the rim.  Accordingly,
\(d_{\rm ext}=T_{\rm ext}^0(0)\) is evaluated with that pin retained.
Prescribed relative guard data are affine equalities, so
Proposition~\ref{prop:comparisons}, with the independent increments
on cut bonds, still gives \(0<T^\tau/d_T\le1\).  To prove the centered
bound in \eqref{geo:kernel-bounds}, impose the flat end event \([H_i]\)
in each incident annulus and choose its representative to be zero on
that end.  The restored-bond factor now depends on the relative integer
translations and the exterior heights, and is independent of the
guard data.  Adding these end constraints decreases the data-to-zero
ratio, giving the product of the annular flat-pin ratios in
\eqref{geo:kernel-bounds}.  This proof uses no unpinned exterior law.
For multiplication, place the exterior, with its full pin constraint,
at the root of the tree of cut blocks.  Its configurations are summed
subject to every retained pin equality.  Once its additive constant is
fixed, the remaining block constants are in bijection with the integer
differences along the tree edges.  Conditioning on the full annular guards
separates their two sides, so the tree multiplication of
\cite[Lemma~5.2]{S} and the identity
\eqref{geo:kernel-sector-identity} follow with this pin retained.

To finish the attachment calculation, let \(e_\tau\) be the expectation
of \(T^\tau_{\rm ext}/d_{\rm ext}\) under the product of the propagated
last-guard laws of the whole nests, and define \(e_0\) in the same way.
Their densities relative to the independent cut laws are at most
\(2^k\), so \eqref{geo:attachment-condition} implies
\(1-e_\tau,1-e_0\le 2^k(1-e^{-2\delta-o(1)})\).
Exact multiplication and cancellation of \(d_{\rm ext}\) give
\begin{equation}\label{geo:attachment-factorization}
 \Phi_{\rm full}^\tau
       =\mathcal N_J^k\Phi_{\rm ext}^\tau\frac{e_\tau}{e_0},
 \qquad
 \log\Phi_{\rm full}^\tau
       =k\log\mathcal N_J+\log\Phi_{\rm ext}^\tau
                              +O_k(\delta)+o(1).
\end{equation}
Here \(\delta\) is small at fixed \(k\), and \(J\to\infty\) with
the core, \(P,\lambda,\varrho_0\) fixed.  The first identity is exact;
the second is the attachment consequence we shall use.

\subsection{A pinned comparison for the thermodynamic correlation}

We first place the desired thermodynamic correlation between two finite
height ratios.  For a large dyadic \(R_{\rm out}\ge8\), let
\[
 \mathsf F_{R_{\rm out}}
       =\overline{Q_{7R_{\rm out}/8}}\setminus Q_{3R_{\rm out}/4}.
\]
The region \(\mathsf F_{R_{\rm out}}\) is a positive-width rim, expressible as an
aligned union of positive-area boxes with dyadic vertices.  In lattice units let
\(G_{R_{\rm out},r}\) be the dual primary array in \(rQ_{R_{\rm out}}\), with coherent
rounding, and let \(\mathsf F_{R_{\rm out},r}\) be the filled primary boxes
representing \(r\mathsf F_{R_{\rm out}}\).  Put
\[
 \Phi_{R_{\rm out},r}^{\rm fr}=\Phi_{G_{R_{\rm out},r}}^\tau,
 \qquad
 \Phi_{R_{\rm out},r}^{\rm pin}
     =\Phi_{G_{R_{\rm out},r},\,h_{\mathsf F_{R_{\rm out},r}}=0}^\tau,
\]
where the sources are \(0,re_1\), with periods \(+2\pi,-2\pi\).

There is one connection suitable for both comparisons.  In the primal
lattice take the actual nearest-neighbor path
\[
 (0,0),(1,0),\ldots,(r,0)
\]
and give its edges the unit current directed from one source to the
other.  The dual connection is supported on the dual bonds crossing
these path edges.  They lie inside \(rQ_{R_{\rm out}/2}\), for all large \(r\),
so the connection vanishes on every bond meeting the rim or lying
outside it.  Use this connection in every larger free spin box.

For a box \(\Lambda_R\) containing \(rQ_{R_{\rm out}}\) in its interior, release
the matching equalities for all bonds crossing the boundary of
\(G_{R_{\rm out},r}\), retaining their independent Bessel increments as in the
finite comparison.  The resulting inner height component has a free
wall; the connection on the exterior component is zero and its ratio
cancels.  Proposition~\ref{prop:comparisons} therefore gives the upper
bound by \(\Phi_{R_{\rm out},r}^{\rm fr}\).  For the lower bound add the
common hard pin \(h_{\mathsf F_{R_{\rm out},r}}=0\) before making any cut.  Added precision
decreases the sector ratio.  This filled rim separates the interior
from the entire exterior.  Because the connection is supported on the
specified path inside the rim, the partition factor for all exterior
variables, including the original fixed exterior face, is identical in
the numerator and denominator.  It cancels.  The same cancellation of
the outer collar occurs in \(G_{R_{\rm out},r}\) with this rim pinned,
so the remaining ratio is exactly \(\Phi_{R_{\rm out},r}^{\rm pin}\).
Hence, for every sufficiently
large \(R\) at each fixed \(R_{\rm out},r\),
\[
 \Phi_{R_{\rm out},r}^{\rm pin}\le
 \E^{\rm free}_{R,b_c}\cos(\theta_0-\theta_{re_1})
 \le \Phi_{R_{\rm out},r}^{\rm fr}.
\]
Taking \(R\to\infty\) at this point proves
\begin{equation}\label{geo:thermodynamic-bracket}
 \Phi_{R_{\rm out},r}^{\rm pin}\le C_{b_c}(r)
                                    \le \Phi_{R_{\rm out},r}^{\rm fr}.
\end{equation}
The finite comparisons just used are the affine and matching-constraint
comparisons of \cite[Lemma~2.2]{C}; the exterior cancellation is why the
location of the connection matters.

Figure~\ref{geo:pinned-schematic} shows the separation used for the
pinned sector estimate.  In the later annular application, its two holes
are those left in the exterior when the two insertion nests are removed
from the full pinned comparator \(\Phi_{R_{\rm out},r}^{\rm pin}\).
Its retained arc \(\Gamma\) is the portion
of the primal source path between the two hole walls.  The guard
separates the strip bonds from the artificial side boundary of the patch,
including at the hole exits; the rim lies beyond the patch.

\begin{figure}[tb]
\centering
\begin{tikzpicture}[font=\small,line cap=round,line join=round]
  \begin{scope}
    \draw[thick] (-1.8,-1.8) rectangle (1.8,1.8);
    \path[fill=gray!25,even odd rule]
      (-1.56,-1.56) rectangle (1.56,1.56)
      (-1.34,-1.34) rectangle (1.34,1.34);
    \draw[gray!65] (-1.56,-1.56) rectangle (1.56,1.56);
    \draw[gray!65] (-1.34,-1.34) rectangle (1.34,1.34);
    \node at (0,1.45) {\(\mathsf F\)};
    \path[fill=blue!5] (-.34,-.12) rectangle (.34,.12);
    \foreach \x in {-.58,.58} {
      \begin{scope}[xshift=\x cm]
        \path[draw,fill=white,thick]
          (-.14,-.28)--(.14,-.28)--(.14,-.18)--(.24,-.18)
          --(.24,.18)--(.14,.18)--(.14,.28)--(-.14,.28)
          --(-.14,.18)--(-.24,.18)--(-.24,-.18)--(-.14,-.18)--cycle;
      \end{scope}
    }
    \draw[blue!65!black,dashed,thick]
      (-.34,.12)--(.34,.12) (-.34,-.12)--(.34,-.12);
    \draw[black!60,densely dotted,line width=.65pt] (-.58,0)--(.58,0);
    \draw[red!70!black,densely dashed,thick] (-.34,0)--(.34,0);
    \fill (-.58,0) circle (1.2pt) (.58,0) circle (1.2pt);
    \node[above] at (-.58,.31) {\(+\)};
    \node[above] at (.58,.31) {\(-\)};
    \node at (0,.78) {\(W\)};
    \draw[-{Stealth[length=1.5mm]},blue!65!black] (0,.65)--(0,.16);
    \draw[black!60,densely dotted,line width=.65pt] (-1.08,-.72)--(-.70,-.72);
    \node[anchor=west,font=\scriptsize] at (-.62,-.72) {full path};
    \draw[red!70!black,densely dashed,thick] (-1.08,-.98)--(-.70,-.98);
    \node[anchor=west,font=\scriptsize] at (-.62,-.98) {retained \(\Gamma\)};
    \node[above] at (0,1.92) {comparison square \(Q_{R_{\rm out}}\)};
    \node[below] at (0,-1.94) {(a) Global arrangement};
  \end{scope}

  \begin{scope}[xshift=4.3cm]
    \path[fill=gray!12] (-.78,-1.55) rectangle (0,1.55);
    \path[fill=blue!5] (0,-1.10) rectangle (3.50,1.10);
    \path[fill=red!7] (0,-.27) rectangle (3.50,.27);

    \path[fill=yellow!38]
      (0,.15) rectangle (3.50,.27) (0,.27) rectangle (3.50,.39)
      (0,-.39) rectangle (3.50,-.27) (0,-.27) rectangle (3.50,-.15);
    \path[fill=blue!25]
      (0,.54) rectangle (3.50,.80) (0,-.80) rectangle (3.50,-.54);
    \path[fill=green!23]
      (0,.98) rectangle (3.50,1.10) (0,1.10) rectangle (3.50,1.22)
      (0,-1.22) rectangle (3.50,-1.10) (0,-1.10) rectangle (3.50,-.98);

    \draw[blue!65!black,dashed,thick]
      (0,1.10)--(3.50,1.10) (0,-1.10)--(3.50,-1.10);
    \draw[red!70!black,thick]
      (0,.27)--(3.50,.27) (0,-.27)--(3.50,-.27);
    \draw[black!60,densely dotted,line width=.65pt] (-.78,0)--(0,0);
    \draw[red!70!black,densely dashed,thick] (0,0)--(3.50,0);

    \fill[black!65] (1.55,.27) circle (.8pt);
    \draw[black!65,line width=.55pt,-{Stealth[length=1.4mm]}]
      (1.55,.27) .. controls (.85,.30) and (.13,.34) .. (.13,.48)
      -- (.13,.85) .. controls (.13,1.12) and (1.08,1.37) .. (1.94,1.37);
    \node[anchor=west,font=\scriptsize,text=black!70] at (2.00,1.37)
      {outside \(W\)};

    \draw[black,line width=.8pt]
      (-.78,1.55)--(0,1.55)--(0,-1.55)--(-.78,-1.55);
    \node[font=\scriptsize,text=black!70] at (-.42,.58) {hole};
    \node[font=\scriptsize,rotate=90,text=black!70] at (-.43,-.69) {wall};

    \draw[black!55,thin]
      (-.10,1.22)--(-.10,1.55) (-.16,1.22)--(-.04,1.22)
      (-.16,1.55)--(-.04,1.55)
      (-.10,-1.55)--(-.10,-1.22) (-.16,-1.55)--(-.04,-1.55)
      (-.16,-1.22)--(-.04,-1.22);
    \node[font=\scriptsize,anchor=east,text=black!65] at (-.18,1.385) {\(>0\)};
    \node[font=\scriptsize,anchor=east,text=black!65] at (-.18,-1.385) {\(>0\)};

    \node[anchor=west] at (3.60,1.10) {\(H'\)};
    \node[anchor=west] at (3.60,.67) {\(\mathsf P\)};
    \node[anchor=west] at (3.60,.27) {\(H\)};
    \draw[red!70!black] (3.50,0)--(3.66,-.20);
    \node[anchor=west] at (3.66,-.23) {\(\Gamma\)};
    \node[anchor=west,font=\scriptsize,text=blue!65!black,align=left] at (3.60,-1.10)
      {artificial side\\boundary of \(W\)};
    \node[font=\scriptsize] at (1.80,-.67) {full data \(h_{\mathsf P}=u\)};
    \node[font=\scriptsize,anchor=east,text=black!65] at (3.50,-1.43)
      {to the other hole \(\longrightarrow\)};
    \node[above] at (1.70,1.92) {plan view at one hole exit};
    \node[below] at (1.70,-1.94) {(b) The guard meets the wall};
  \end{scope}
\end{tikzpicture}
\caption{Schematic geometry for the pinned comparison.  The dotted primal
path includes the portions inside the holes; its red portion \(\Gamma\) is
the retained arc between the walls.  At the enlarged exit, both guard bands
\(\mathsf P\) reach the original wall between \(H\) and \(H'\); conditioning
fixes all primary guard heights \(h_{\mathsf P}=u\).  A retained route from
a strip interface to the artificial side boundary of \(W\) must therefore
meet a guard.  The red interfaces lie in \(H\), and \(H'\) straddles that
dashed boundary.  Endpoint boxes are separate across each cut and lie in
their respective components before pins are imposed.  The buffers from \(H\) to \(\mathsf P\),
from \(\mathsf P\) to \(H'\), and at side corners are positive.  The same
order continues around bends to the other wall, and the rim \(\mathsf F\)
lies beyond \(W\).  Distances are schematic.}
\label{geo:pinned-schematic}
\end{figure}

\begin{samepage}
To estimate the lower comparator, we need a sector lower bound with
the rim pinned.  The following lemma supplies exactly that bound.

\begin{lemma}[Trial energy with a pinned rim]\label{geo:pinned-sector}
Let \(U\) be a bounded connected polygonal Lipschitz domain in the
fixed-geometry class of Proposition~\ref{prop:height-inputs}.  Its outer
boundary and the boundaries of its two holes are simple and rectilinear,
with dyadic vertices and positive side lengths; the hole closures are
disjoint and strictly inside the outer boundary, leaving positive-width
passages.  Let \(\mathsf F\) be an aligned union of positive-area boxes
with dyadic vertices in \(U\), permitted as a pin in that fixed geometry.
Suppose a simple dyadic rectilinear arc joins
interiors of sides of the two holes transversely and has a neighborhood
at positive distance from \(\mathsf F\) and all nonincident boundaries.
Give the holes periods \(+2\pi,-2\pi\).
Let \(\vartheta=dv\) be a smooth closed one-form near \(\overline U\)
with these periods, with a multivalued primitive whose only jump is
across that arc and whose branch is zero on a neighborhood of
\(\mathsf F\).  Choose the canonical integral connections obtained
from the jumps of this branch along a coherent rounded discretization
of the arc, or connections
related to them by integral gauges whose potentials vanish on every
component of \(\mathsf F\).  Then, along every admissible rounded mesh sequence,
\begin{equation}\label{geo:pinned-sector-bound}
 \liminf_{\nu\to\infty}
       \Phi_{\mathcal G_\nu(U),\,h_{\mathsf F}=0}^\tau
       \ge \exp\left\{-\frac1{2a}\int_U|\vartheta|^2\right\}.
\end{equation}
\end{lemma}
\end{samepage}

\begin{proof}
We adapt the confined-sector argument in the proof of
\cite[Proposition~3.4]{S}, retaining the rim pin.  A permitted gauge
preserves the pin and the ratio, so we use the canonical connection.
Round the arc within a bounded number of lattice steps and extend its
transverse ends to the rounded hole walls.  Transport the local branches
across this rounding, which remains away from \(\mathsf F\), and sample
them at the lattice vertices,
using a coherent one-sided assignment when the rounded cut meets a
vertex.  Choose the integral jump \(\tau_e\) on each edge with its
matching sign.  The resulting real lattice values \(v_\nu\) satisfy
\[
 \nabla_e v_\nu+\tau_e=\int_e\vartheta,
 \qquad v_\nu=0\text{ on the rounded pin }\mathsf F.
\]
The identity follows by integrating \(dv\) along each edge, including
its integral jump when the edge crosses or meets the rounded cut.
The branches differ by integral multiples of \(2\pi\), so
\(\tau_e\in2\pi\Z\); supported bonds lie within one lattice step of
the rounded arc.  The chosen branch is
zero near every component of the pin.  These branches are fixed before
refining the lattice.

Isolate a narrow rectilinear strip following the arc between the holes
by cutting the bonds across both of its sides.  Choose a wider strip
patch \(W\) and full guard bands \(\mathsf P\) on the two sides of the
narrow strip, with positive buffers from the restored bonds and from
the artificial boundary of \(W\).  The strip and patch meet the hole
walls in interiors of sides.  Choose all their widths and coordinates
dyadically and keep \(W\) at positive distance from \(\mathsf F\).
For every sufficiently fine mesh, the narrow strip contains all bonds
supporting \(\tau\).  After this cut the strip admits a single-valued
lift.  The remaining domain excludes that entire support, so its lift
is the original branch, already zero near every component of
\(\mathsf F\).  The patch \(W\), which is disjoint from the pin, admits
a lift both before and after restoring the strip.
Fix an endpoint region \(H\) inside \(W\) containing both primary
endpoints of every restored strip bond, and a region \(H'\) containing
both endpoints of every bond across the artificial boundary of \(W\).
Choose these as aligned unions of positive-area boxes with dyadic vertices, separated from
\(\mathsf P\).  At the transverse exits, their intervals on the hole
sides follow the same nested order as the strip, guard, and patch
boundaries, with positive gaps from each other and the side corners.
The guard bands continue to the original hole walls between the
intervals for \(H\) and \(H'\).  Hence every path in the retained
domain from a restored strip interface to the artificial boundary
of \(W\) meets a guard, including paths near the exits.
Use separate boxes on the two sides of every pre-pin cut, so each box
lies in one pre-pin component.  The clearance from \(W\) to the rim
also permits \(H'\) to be disjoint from \(\mathsf F\).

Fix one dyadic cell side \(\ell>0\) and one cell partition aligned with
the holes, the rim, the strip, \(W\), and all guard and endpoint bands.
At each topology determine its components before imposing pins.  The
restriction of this same partition has every whole observation cell in
one such component.  Using one box average on every cell confines the
constant of every pre-pin component, including those later fixed by
\(\mathsf F\).  On every topology, with \(\mathsf F\)
already pinned, add \(\frac12\sum_Q(h_Q-(v_\nu)_Q)^2\) in the numerator and
\(\frac12\sum_Qh_Q^2\) in the denominator.  Denote the resulting full
graph ratio by \(\Phi^{\tau,\mathrm{conf}}_{\nu,\ell}\); below, a topology
subscript specifies its cut version.  The new observation
residuals have the same positive precision, and the pin is compatible
since \(v_\nu=0\) there.  Thus
\begin{equation}\label{geo:confinement-below-bare}
             0<\Phi^{\tau,\mathrm{conf}}_{\nu,\ell}
                 \le \Phi_{\mathcal G_\nu(U),\,h_{\mathsf F}=0}^\tau.
\end{equation}
On a cut component \(V\), choose its single-valued lift \(v_V\) and
sample it at the lattice vertices with \(v_V-v_\nu\in2\pi\Z\)
pointwise.  Since \(\nabla_e v_V=\int_e\vartheta\), the integral
potential \(\psi=v_\nu-v_V\) satisfies \(\tau+\nabla\psi=0\).
If \(V\) meets \(\mathsf F\), both lifts use the zero branch there,
so \(\psi=0\) on the pin.  Writing \(B\) for the cell-average map,
the gauge rule of Section~\ref{sec:heights} changes the observation
center \(Bv_\nu\) to \(Bv_\nu-B\psi=Bv_V\).  The mixed height limit
in Proposition~\ref{prop:height-inputs}, at this fixed \(\ell\),
and Gaussian completion give for its confined ratio the limit
\(\exp\{-I_{\ell,V}(v_V)/2\}\), where
\[
 I_{\ell,V}(v_V)=
 \inf_{\substack{u\in H^1(V)\\u=0\ \mathrm{on}\ V\cap\mathsf F}}
 \left\{a^{-1}\int_V|\nabla u|^2+
       \sum_{Q\subset V}(u_Q-(v_V)_Q)^2\right\}.
\]
An empty pin in this display imposes no constraint.  The unpinned
constant is integrated.  The converging lattice cell centers can be used
in the Laplace limit: transforms at nearby fixed arguments give local
equicontinuity by H\"older's inequality.  Trying \(u=v_V\) is allowed,
and cell Poincar\'e gives, for bounded values of the displayed functional,
\[
 \|u-v_V\|_{L^2(V)}^2
 \le C\ell^2\left(\|\nabla u\|_2^2+\|\nabla v_V\|_2^2+
                    \sum_Q(u_Q-(v_V)_Q)^2\right).
\]
Weak compactness and lower semicontinuity therefore show that
\(I_{\ell,V}(v_V)\to a^{-1}\int_V|\nabla v_V|^2\) as
\(\ell\downarrow0\).  This is the single-lift argument using the
hard-pin energy space of \cite[Theorem~3.1]{C}.

We verify that restoring the strip has a vanishing averaged loss.
For full absolute data \(u\) on \(\mathsf P\), let
\[
 K^\tau(u)=\frac{Z_{\rm open,\mathsf F}^\tau(u)}
                 {Z_{\rm cut,\mathsf F}^\tau(u)},
 \qquad d_0=K^0(0),
\]
with the confinement retained.  Here \(\mathrm{open}\) restores the bonds
cut across the two sides of the narrow strip, while \(\mathrm{cut}\)
retains those cuts.  The guard separates the restored bonds
from the artificial side boundary of \(W\) and from \(\mathsf F\).
No observation cell crosses a guard division.  Conditioning on the full
guard hence makes
\(K^\tau(u)\) exactly the kernel computed on \(W\) alone, so it is
unchanged by the exterior pin.  Finite affine comparison gives
\(0<K^\tau(u)\le d_0\).

Adding the pin on the endpoint region \(H\) to the centered
data-to-zero comparison makes the restored factor a scalar.
Bayes' formula then gives
\[
 1\ge\E_{{\rm cut}\,W,0}\frac{K^0}{d_0}
       \ge e^{-D_{\nu,\ell}^{W}(H,\mathsf P)}.
\]
Here \(D_{\nu,\ell}^{W}\) is the absolute pin interaction in the centered
cut patch with its component-local cell observations, before either
hard pin is imposed.  The localization estimate of
\cite[Lemma~3.2]{S} applies to this pre-pin law with no existing hard
pins: for every fixed dyadic cut mixed geometry and positively separated
aligned pin regions that are unions of positive-area boxes with dyadic vertices,
\begin{equation}\label{geo:cell-pin-localization}
 \lim_{\ell\downarrow0}\limsup_{\nu\to\infty}
                  D_{\nu,\ell}(H,\mathsf P)=0.
\end{equation}
It uses the Gaussian pin limit of \cite[Theorem~3.2]{C} for precision
\(a^{-1}\int|\nabla u|^2+\sum_Q u_Q^2\); \(\ell\) decreases through
aligned dyadic values after the mesh limit.  Both topologies of \(W\)
have the single-lift limit just proved, with the same energy since the
deleted interfaces have zero area.  Consequently the nonnegative patch
loss satisfies
\begin{align*}
 L_W^\tau
 &=\sum_u\frac{Z_{{\rm cut}\,W}^\tau(u)}{Z_{{\rm cut}\,W}^0}
                    \left(1-\frac{K^\tau(u)}{d_0}\right)\\
 &=\Phi_{{\rm cut}\,W}^{\tau,\mathrm{conf}}-
       \Phi_{{\rm open}\,W}^{\tau,\mathrm{conf}}\,
                    \E_{{\rm cut}\,W,0}\frac{K^0}{d_0}
       \longrightarrow0
       \qquad(\nu\to\infty,\text{ then }\ell\downarrow0).
\end{align*}
The centered loss has the same conclusion.

It remains to pass this loss from \(W\) to the whole pinned graph.
Let \(G\) be the current graph with the narrow strip cut, and let
\(G'\) further cut the bonds across the artificial boundary of \(W\).
At prescribed guard data, finite comparison, with the rim pin retained,
gives
\begin{equation}\label{geo:global-data-comparison}
 \frac{Z_{G,\mathsf F}^\tau(u)}{Z_{G,\mathsf F}^0}
 \le
 \frac{\Prob_{G,\mathsf F,0}(h_{\mathsf P}=0)}
      {\Prob_{G',\mathsf F,0}(h_{\mathsf P}=0)}
 \frac{Z_{G',\mathsf F}^\tau(u)}{Z_{G',\mathsf F}^0}.
\end{equation}
Use the endpoint region \(H'\) already chosen around that artificial
boundary, separated from \(\mathsf P\).
If \(R(u)\) denotes the centered restoration factor from \(G'\) to
\(G\), summation gives the probability ratio in
\eqref{geo:global-data-comparison} as \(R(0)/\E_{G',\mathsf F,0}R\).
The finite comparison argument in the proof of \cite[Lemma~3.9]{C}
can be applied with the common rim constraint retained: adding the
pin on \(H'\) fixes both endpoints of each restored bond, so its
restoration factor becomes independent of the guard data.  Dividing
the data-to-zero ratios and averaging by Bayes' formula yields
\[
 \frac{\E_{G',\mathsf F,0}R}{R(0)}
       \ge e^{-D^{\mathsf F}_{\nu,\ell}(H',\mathsf P)}.
\]
To control this conditional interaction by the available unpinned
theorem, use the centered confined law on the pre-pin graph \(G'\) and
write all pin events with zero values.  Bayes' formula gives exactly
\begin{equation}\label{geo:conditional-pin-identity}
 \begin{split}
 0\le D^{\mathsf F}_{\nu,\ell}(H',\mathsf P)
 &:=\log\frac{\Prob(h_{H'}=0\mid h_{\mathsf P}=0,h_{\mathsf F}=0)}
                 {\Prob(h_{H'}=0\mid h_{\mathsf F}=0)}\\
 &=D_{\nu,\ell}(H'\cup\mathsf F,\mathsf P)
                   -D_{\nu,\ell}(\mathsf F,\mathsf P)\\
 &\le D_{\nu,\ell}(H'\cup\mathsf F,\mathsf P).
 \end{split}
\end{equation}
The first inequality is centered pin monotonicity in the law already
pinned on \(\mathsf F\); the last uses the nonnegativity of the
unconditional interaction.  The union \(H'\cup\mathsf F\) remains
at positive distance from \(\mathsf P\).  Its boxes and the observation
cells are component-local on \(G'\) before pins are imposed.
Thus \eqref{geo:cell-pin-localization}, applied to this union in that
pre-pin geometry, makes the last expression tend to zero in the stated
order.  This derives the needed conditional estimate from the
unconditional pin theorem.

Multiply \eqref{geo:global-data-comparison} by
\(1-K^\tau(u)/d_0\ge0\) and sum over \(u\).  On \(G'\), the part outside
\(W\) factors off as a confined shifted-to-centered ratio, at most one
by Proposition~\ref{prop:comparisons}.  The global loss is therefore
at most \(e^{D_{\nu,\ell}(H'\cup\mathsf F,\mathsf P)}L_W^\tau\), and tends
to zero.  The same argument applies in the centered sector.  If these
two losses, normalized by the centered cut sum, are \(L_\tau,L_0\),
exact summation of the restored sums gives
\[
              \Phi_{\rm open}^{\tau,\mathrm{conf}}
                 =\frac{\Phi_{\rm cut}^{\tau,\mathrm{conf}}-L_\tau}
                        {1-L_0}.
\]
The positive scalar \(d_0\) has canceled.  The losses vanish, and the
cut components' single-lift energies add to \(\int_U|\vartheta|^2\).
This proves the ordered estimate
\[
 \lim_{\ell\downarrow0}\limsup_{\nu\to\infty}
 \left|\Phi^{\tau,\mathrm{conf}}_{\nu,\ell}
 -\exp\left\{-\frac1{2a}\int_U|\vartheta|^2\right\}\right|=0.
\]
For each fixed aligned \(\ell\),
\eqref{geo:confinement-below-bare} bounds the bare pinned ratio below
by this confined ratio.  Taking its lattice liminf and then letting
\(\ell\downarrow0\) proves \eqref{geo:pinned-sector-bound}.
The holes, rim, strip, patch, and bands were all fixed before these limits.
\end{proof}

\subsection{The energies outside the nests}

Let \(D=(-1,1)^2\) be the reference square, and write its Dirichlet
Green kernel as
\[
 G_D(z,w)=\frac1{2\pi}\log\frac1{|z-w|}+R_D(z,w).
\]
The regular part \(R_D\) is smooth across the diagonal on compact
subsets of \(D\times D\).  This Dirichlet kernel is the stream function
for the free-height variational problem: rotating the gradient of
\(2\pi G_D\) gives the unit-period trial field tangent to the outer
wall.  The exterior-energy calculation of \cite[Lemma~6.1]{S} gives,
for one circular hole of radius \(t\) at the center,
\begin{equation}\label{geo:single-energy}
 \mathcal M_{\rm cen}(t)
       =2\pi\log(1/t)+4\pi^2R_D(0,0)+O(t).
\end{equation}
More generally, for separated unit charges \(\sigma_i\) at \(z_i\),
that lemma gives the constant term
\(4\pi^2[\sum_iR_D(z_i,z_i)+2\sum_{i<j}\sigma_i\sigma_jG_D(z_i,z_j)]\)
and the leading term \(2\pi k\log(1/t)\).  The error is \(O(t)\) at
each fixed separated configuration.  It follows by integrating the
Green energy on the small circles; the regular harmonic part has zero
total normal derivative there, so no additional logarithmic error is
created.

For the pair \(0,e_1\) in \(Q_{R_{\rm out}}=R_{\rm out}D\), put
\[
 \mathfrak h_{R_{\rm out}}
   =R_{Q_{R_{\rm out}}}(0,0)+R_{Q_{R_{\rm out}}}(e_1,e_1)
                           -2G_{Q_{R_{\rm out}}}(0,e_1).
\]
Green scaling gives
\(R_{Q_{R_{\rm out}}}(z,w)=(2\pi)^{-1}\log R_{\rm out}
                         +R_D(z/R_{\rm out},w/R_{\rm out})\).
The singular logarithm between \(0\) and \(e_1\) is zero, so
\[
 \mathfrak h_{R_{\rm out}}
   =R_D(0,0)+R_D(e_1/R_{\rm out},e_1/R_{\rm out})
                           -2R_D(0,e_1/R_{\rm out})
                              \longrightarrow0.
\]
The \(\log R_{\rm out}\) terms have canceled by neutrality.  The same exterior
formula therefore reads
\begin{equation}\label{geo:free-pair-energy}
 \mathcal M_{{\rm pair},R_{\rm out}}(t)
        =4\pi\log(1/t)+4\pi^2\mathfrak h_{R_{\rm out}}
                                                +O_{R_{\rm out}}(t).
\end{equation}
In both \eqref{geo:single-energy} and \eqref{geo:free-pair-energy},
replacing circular holes by \(tP\) changes the right side by
\(O(\varepsilon)\): restrict admissible forms between the exteriors
of the circles of radii \(t(1-\varepsilon)\) and
\(t(1+\varepsilon)\), and use their logarithmic terms.  The constant
in this error is independent of small \(t\).  The same restriction
argument on a circular annulus proves, for either unit period,
\begin{equation}\label{geo:annular-energy}
 \mathcal M_{\lambda P\setminus P^\circ}(\pm1)
                   =2\pi\log\lambda+O(\varepsilon),
\end{equation}
with an error independent of large fixed \(\lambda\).  This is also
the annular energy estimate in the proof of
\cite[Proposition~5.3]{S}.

For the pinned lower bound we exhibit a trial attaining the same pair
asymptotics.  Identify the plane with \(\mathbb C\), and let
\[
                  \Theta(z)=\arg z-\arg(z-1).
\]
Use the branch on \(\mathbb C\setminus[0,1]\) that tends to zero at
infinity.  Its gradient has periods \(+2\pi,-2\pi\), and
\(\Theta(z)=O(|z|^{-1})\), \(\nabla\Theta(z)=O(|z|^{-2})\) at
infinity.  Rotation identifies \(\nabla\Theta\), up to orientation,
with the gradient of \(\log|z|-\log|z-1|\).  Green's identity on the
plane outside the two radius-\(t\) disks gives
\begin{equation}\label{geo:plane-pair-energy}
 \int_{\R^2\setminus(B(0,t)\cup B(e_1,t))}|\nabla\Theta|^2
                     =4\pi\log(1/t)+O(t).
\end{equation}
Indeed each small circle contributes \(2\pi\log(1/t)+O(t)\).
The cross logarithm is \(\log|e_1|=0\), and the boundary term at
infinity vanishes by the displayed decay.  Harmonicity of each regular
part removes a logarithmically amplified error, as in the preceding
Green calculation.

Choose a smooth cutoff \(\chi_{R_{\rm out}}\) equal to one on
\(Q_{R_{\rm out}/2}\), zero outside \(Q_{2R_{\rm out}/3}\), and satisfying
\(|\nabla\chi_{R_{\rm out}}|\le C/R_{\rm out}\).
The cut \([0,1]\) lies where \(\chi_{R_{\rm out}}=1\).  Thus
\(\vartheta_{R_{\rm out}}=d(\chi_{R_{\rm out}}\Theta)\) is a well-defined smooth closed
one-form off the holes: its two branches have the same differential
across that cut.  It has the required periods and a lift identically
zero near \(\mathsf F_{R_{\rm out}}\).  The energy of both the discarded tail and
the cutoff correction is \(O(R_{\rm out}^{-2})\), since on the transition
region \(|\Theta\nabla\chi_{R_{\rm out}}|+|\nabla\Theta|
\le C R_{\rm out}^{-2}\) and its area is \(O(R_{\rm out}^2)\).  The polygon inclusions
\eqref{geo:polygon} and \eqref{geo:plane-pair-energy} consequently give
\begin{equation}\label{geo:pinned-trial-energy}
 \int_{Q_{R_{\rm out}}\setminus(tP\cup(e_1+tP))}|\vartheta_{R_{\rm out}}|^2
       \le4\pi\log(1/t)+C\bigl(t+\varepsilon+R_{\rm out}^{-2}\bigr).
\end{equation}
The constants here are uniform for large \(R_{\rm out}\) and sufficiently small
\(t,\varepsilon\).

We record carefully how these energies apply when the last radius moves
with the lattice.  Suppose \(t_\nu\to t>0\).  Around each limiting source,
choose fixed dyadic stair holes lying just inside the disk of radius
\((1-3\varepsilon)t\) and just outside the disk of radius
\((1+3\varepsilon)t\).  Small further perturbations of their dyadic
centers and sides give strict eventual inclusions of the rounded moving
holes \(t_\nu P\).  For the pair, choose exit sides whose interiors
meet the fixed source-to-source path transversely.
All these holes contain the source, remain separated,
and leave positive passages.  Cut comparison brackets the three sector
ratios, using restrictions of the one connection chosen before the cuts.
For the pinned ratios keep exactly the same rim in both brackets;
the rim is connected, so the integral potential relating two flat
representatives with the same periods and zero connection near the rim
is constant there and can be normalized to zero.  Apply
\eqref{geo:free-sector} to the fixed free exteriors and
Lemma~\ref{geo:pinned-sector} to the fixed pinned inner bracket.
The changes of radius cost \(O(\varepsilon)\) in the energy formulas.
This is the moving-hole argument of \cite[Lemma~6.2]{S}, with the pin
retained in the lower comparison.  Bounded rounding of the outer arrays
and filled rim boxes is included in the fixed mixed limits.  In
particular, every sector limit just invoked has holes of positive radius.

Let \(\Phi_{{\rm cen},\nu}^{\rm ext}(t_\nu)\) denote the free exterior in
the reference square \(D\), with one moving polygonal hole at the center.
The preceding argument gives
\begin{equation}\label{geo:center-exterior}
 \limsup_{\nu\to\infty}\left|
 \log\Phi_{{\rm cen},\nu}^{\rm ext}(t_\nu)+
       \frac{2\pi\log(1/t)+4\pi^2R_D(0,0)}{2a}\right|
                      \le C(\varepsilon+t).
\end{equation}
For the pair exterior in \(Q_{R_{\rm out}}\), write
\(\Phi_{{\rm pair},R_{\rm out},\nu}^{\rm ext,fr}(t_\nu)\) for the free-wall ratio and
\(\Phi_{{\rm pair},R_{\rm out},\nu}^{\rm ext,pin}(t_\nu)\) for the ratio with
\(\mathsf F_{R_{\rm out}}\) pinned.  Releasing that pin bounds the latter by the
former.  Combining \eqref{geo:free-pair-energy},
\eqref{geo:pinned-trial-energy}, and Lemma~\ref{geo:pinned-sector}
shows the following useful uniform statement.  Given \(\delta>0\),
first choose \(R_{\rm out}\) sufficiently large, then choose a bound
\(\eta>0\) and \(\varepsilon>0\) sufficiently small.  For either
boundary choice \(\star\in\{\mathrm{fr},\mathrm{pin}\}\) and every
sequence \(t_\nu\to t\in(0,\eta]\),
\begin{equation}\label{geo:pair-exterior}
 \limsup_{\nu\to\infty}\left|
    \log\Phi_{{\rm pair},R_{\rm out},\nu}^{\mathrm{ext},\star}(t_\nu)
                      +\frac{4\pi\log(1/t)}{2a}\right|\le\delta.
\end{equation}
Here one uses \(\mathfrak h_{R_{\rm out}}\to0\) and \(R_{\rm out}^{-2}\to0\) first, then the
\(O_{R_{\rm out}}(t)\) error.  The bound follows for the pinned ratio from its
lower trial bound and its free upper bound; it asserts no separate
pinned sector theorem.  It is uniform over the stated positive values
of \(t\) because the energy errors are bounded by
\(C_{R_{\rm out}}\eta+C\varepsilon\) and \(C(\eta+\varepsilon+R_{\rm out}^{-2})\).
For each such value the fixed stair brackets are chosen after its
limit has been extracted.

\subsection{Cancellation of the microscopic costs}

\begin{proof}[Proof of Proposition~\ref{geo:comparison}]
First fix integer sequences \(n,n'\to\infty\) with \(n'/n\to d>0\).
Choose a small macroscopic bound \(\eta\), so nests of outer radius at
most \((1+\varepsilon)\eta n\) fit in both squares for all late terms.
After fixing the annular parameters, put
\[
              J=\max\{i:\varrho_i\le\eta n\}.
\]
Then \(J\to\infty\).  Every subsequence has a further subsequence on
which \(\varrho_J/n\to t\in[\eta/\lambda,\eta]\); in the other
square the same physical nest has scaled radius tending to \(t/d\).
Both limits are positive with these parameters fixed.

We check the attachment condition for either square.  After joining
the last annulus to its exterior, the hole radius changes from \(t\)
to \(t/\lambda\).  By \eqref{geo:free-sector} and
\eqref{geo:annular-energy}, its annular log ratio is
\(-2\pi\log\lambda/(2a)+O(\varepsilon)+o(1)\).  In
\eqref{geo:center-exterior} the two exterior constant terms agree,
and the difference of their radius energies is exactly
\(2\pi\log\lambda\).  Thus the leading terms in
\(\log\mathfrak q^\tau_{\rm ext}\) cancel and its limiting absolute error is
\(O(\varepsilon+\eta)\).  In the other square the error is
\(O_d(\varepsilon+\eta)\), using \(t/d\) and \(t/(d\lambda)\).
These bounds hold uniformly for the possible \(t\) at large fixed
\(\lambda\).  The flat-pin interactions can also be made arbitrarily
small by choosing \(\lambda\) large.  More precisely, for any desired
attachment accuracy, choose \(\varepsilon,\eta\) small, then
\(\lambda\) large for the flat-pin estimate, and finally the fixed
\(\varrho_0\) large for \eqref{geo:regular-loss}.  The inequalities
\eqref{geo:attachment-condition} then hold with arbitrarily small
\(\delta\).

Apply \eqref{geo:attachment-factorization} with one nest to each
square, using the same \(\mathcal N_J\).  Subtraction cancels this
microscopic ratio.  Formula \eqref{geo:center-exterior} cancels its
regular Green term and gives, up to an arbitrarily small limiting error,
\[
 \log M_{n'}-\log M_n
        =-\frac{2\pi}{2a}\log\frac{n'}n+o(1).
\]
The claimed convergence follows because \(\pi/a=1/8\).
For precision, the displayed \(o(1)\) after arbitrary accuracy means
the following: any subsequence with a fixed positive error would,
after choosing that accuracy, have a further subsequence with the
above limit \(t\) and a smaller error.  This is impossible.  The
argument proves \eqref{geo:center-ratio} for the full sequences.

Taking \(n'=2n\) and summing logarithms on powers of two gives
\(\log M_{2^i}=-(i/8)\log2+o(i)\).  The ratios
\(M_n/M_{2^i}\), \(2^i\le n<2^{i+1}\), stay bounded above and below
by positive constants for large \(i\): otherwise a sequence of
\(n/2^i\) has a convergent subsequence in \([1,2]\) contradicting
\eqref{geo:center-ratio}.  This proves \(M_n=n^{-1/8+o(1)}\).

For the bulk comparison, fix an arbitrary small accuracy \(\delta\).
Choose \(R_{\rm out}\) large and then \(\eta,\varepsilon\) small so that
\eqref{geo:pair-exterior} and \eqref{geo:center-exterior} have errors
as small as required in terms of \(\delta\).  Take \(\eta\) small
enough for two disjoint nests about \(0,e_1\) and for the center nest
in \(D\).  Fix \(P\), next choose \(\lambda\) large for the flat-pin
interactions, and finally choose \(\varrho_0\) large for all regular
annuli.  For \(\Phi_{R_{\rm out},r}^{\rm fr}\),
\(\Phi_{R_{\rm out},r}^{\rm pin}\), and \(M_r\), use the identical nests of depth
\(J=\max\{i:\varrho_i\le\eta r\}\).  The connection chosen on the
primal source-to-source path restricts to one unit of the appropriate
sign in each pair nest.  Its bare ratio is therefore precisely the
same \(\mathcal N_J\) as the center nest's ratio.

Along a further subsequence, \(\varrho_J/r\to t\in
[\eta/\lambda,\eta]\).  The pair exterior before and after the
last-annulus restoration has radii \(t\) and \(t/\lambda\), with the
same outer wall and, in the lower comparison, the same rim pin.
Both are covered by \eqref{geo:pair-exterior}.  The two last annuli
are covered by \eqref{geo:annular-energy}.  In their joining ratio
the leading logarithms cancel exactly:
\[
 -\frac{4\pi\log(\lambda/t)}{2a}
 +\frac{4\pi\log(1/t)}{2a}
 +\frac{2(2\pi\log\lambda)}{2a}=0.
\]
Its log is at most zero by cut comparison and at least an arbitrarily
small negative constant by these exterior bounds.  This verifies
\eqref{geo:attachment-condition} for both pair comparators.  The
center attachment was already checked.  The argument with the pinned
root following \eqref{geo:attachment-condition} permits its use for
\(\Phi_{R_{\rm out},r}^{\rm pin}\).

Equations \eqref{geo:attachment-factorization},
\eqref{geo:pair-exterior}, and \eqref{geo:center-exterior} now give,
for either \(\star=\mathrm{fr}\) or \(\star=\mathrm{pin}\),
\begin{align*}
 \log \Phi_{R_{\rm out},r}^{\star}
    &=2\log\mathcal N_J-\frac{4\pi\log(1/t)}{2a}
                                      +O(\delta)+o(1),\\
 2\log M_r
    &=2\log\mathcal N_J-
          \frac{4\pi\log(1/t)+8\pi^2R_D(0,0)}{2a}
                                      +O(\delta)+o(1).
\end{align*}
The microscopic terms and the radius terms cancel.  The same
subsequence argument as above makes these bounds valid as limiting
bounds for the full sequence.  Since every quantity is positive,
\eqref{geo:thermodynamic-bracket} then implies
\[
 \limsup_{r\to\infty}\left|
       \log\frac{C_{b_c}(r)}{M_r^2}-\frac{4\pi^2}{a}R_D(0,0)\right|
       \le C\delta.
\]
The accuracy was arbitrary, so \eqref{geo:bulk-ratio} holds with
\begin{equation}\label{geo:shape-constant}
             c_{\rm shape}=\exp\left\{\frac{4\pi^2}{a}R_D(0,0)\right\}
                            \in(0,\infty).
\end{equation}
Its finiteness follows from regularity of \(R_D\) at the center.
All mesh limits used a fixed \(R_{\rm out},P,\lambda,\varrho_0\) and a positive
subsequential radius \(t\).  The geometric parameters varied only
through the final accuracy argument, after the thermodynamic limit
had already been taken in \eqref{geo:thermodynamic-bracket}.
\end{proof}
